\subsection{为什么七维球面需要别的信息}
称光滑闭流形为怪球（exotic sphere），是指它与标准球面同胚但不微分同胚。“同胚”和“微分同胚”不应只因两者都有球面外形而合并。Milnor于1956年构造了七维这样的例子；这是外部定理，本章只说明Chern--Weil初等检测在此处的局限，不声称完成怪球构造或分类。

若一个七维流形在光滑同伦意义下与$S^7$等价，第9章使其de Rham上同调仅在次数0、7非零。正次Pontryagin类的次数为$4,8,\ldots$，故全部为零；一般Chern--Weil偶次类也无法在次数7给出顶次示性数。对只给同胚的怪球使用同样的上同调结论时，还须采用本书未证明的拓扑/de Rham比较与相关光滑逼近，不把这一比较隐藏在“同胚不变”四字中。

这并不否定Pontryagin理论在Milnor证明中的作用。区别在于，他使用八维填充及相对配对，而不是仅积分七维流形自身的四形式。为准确说明这一点，给定紧致定向八维$W$，$\partial W=M^7$，并假设整上同调$H^3(M;\mathbb Z)=H^4(M;\mathbb Z)=0$。相对上同调长正合列使
\[
 j:H^4(W,M;\mathbb Z)\longrightarrow H^4(W;\mathbb Z)
\]
为同构。\textbf{以下属于外引拓扑工具：}取整Pontryagin类$p_1^{\mathbb Z}(TW)$，定义
\[
 q(W)=\left\langle\bigl(j^{-1}p_1^{\mathbb Z}(TW)\bigr)^2,[W,M]\right\rangle,
 \qquad\lambda(M)=2q(W)-\sigma(W)\pmod7,
\]
其中$\sigma(W)$是相对四维交叉配对的签名。Milnor原论文定理1断言此剩余类与填充的选择无关，并以它判别某些球面束构造所得流形的光滑结构。该论证使用整类、相对基本类、填充存在性和Hirzebruch签名定理；这些并没有由本章实形式的联络无关性自动推出。

因此不能用$H^4(M;\mathbb R)=0$断言$q(W)=0$，也不能把某张卡上的Chern--Simons三形式随意积分后叫作上述$\lambda$。本章不证明一般Chern--Gauss--Bonnet、Euler类与零点指标比较、整类比较、Hirzebruch签名公式、怪球分类或高维手术理论。第14章已经证明的曲面公式和本章算出的球面/射影积分，均不冒充这些整体定理的一般证明。
